\documentclass[10pt]{report}
\usepackage{epsf}
\usepackage{amsmath}
\usepackage{amssymb}
\usepackage{palatino}
\usepackage[pdftex]{graphics}
\usepackage{fancyhdr}
\usepackage{epsfig}
\usepackage{multirow}
\usepackage{multicol}
\usepackage{cancel}
\usepackage{hyperref}
\usepackage{longtable}
\usepackage{ulem}
\usepackage{draftwatermark}
\usepackage{termcal}
\usepackage{etoolbox}
\patchcmd{\endcalendar}{[l]}{[c]}{}{}

\parindent 0in
\parskip 1ex
\oddsidemargin  0in
\evensidemargin 0in
\textheight 8.5in
\textwidth 6.5in
\topmargin -0.25in

\pagestyle{fancy}
\newcommand{\leftheader}{MedTech Prototyping Skills (BME290L)}
\newcommand{\rightheader}{Palmeri (Spring 2023)}
\lhead{\textbf{\leftheader}}
\rhead{\textbf{\rightheader}}
\lfoot{Compilation Date: \today}
\cfoot{}
\rfoot{\thepage}

\newcommand{\Holiday}[2]{%
    \options{#1}{\noclassday}
    \caltext{#1}{#2 (No Class)}
}

\begin{document}

\section*{Class Syllabus}

\subsection*{Personnel}
\begin{itemize}
    \item[] \textbf{Instructor:} Dr. Mark Palmeri
    \begin{itemize}
        \item While you are welcome to email me
        (\href{mailto:mark.palmeri@duke.edu}{mark.palmeri@duke.edu}), my response times will be way
        faster via Teams.
        \item Office Hours: \url{http://meet.palmeri.io}
        \begin{itemize}
            \item 258 Hudson Hall Annex
            \item \url{https://duke.zoom.us/my/mark.palmeri}
        \end{itemize}
    \end{itemize}
    \item[] \textbf{Lab Master:} Matt Brown \href{mailto:matt.brown@duke.edu}{(matt.brown@duke.edu)}
    \item[] \textbf{Teaching Assistants:} 
    \begin{itemize}
        \item Rebecca Hogewood \href{mailto:rebecca.hogewood@duke.edu}{(rebecca.hogewood@duke.edu)}
        \item Nikhita Gopisetty \href{mailto:nikhita.gopisetty@duke.edu}{(nikhita.gopisetty@duke.edu)}  
        \item Nikhil Gadiraju \href{mailto:nikhil.gadiraju@duke.edu}{(nikhil.gadiraju@duke.edu)}
    \end{itemize}
\end{itemize}

\subsection*{Course Times \& Locations}
\begin{itemize}
    \item[] \textbf{Lecture:} Tues \& Thurs from 12:00--13:15 in Wilkinson 126

    \item[] \textbf{Lab:} Thurs from 13:45--16:45 in the POD or Fri from 08:30--11:30 in Chesterfield 5704
    \footnote{A campus shuttle runs to/from the E-quad; details can be found here:
\url{https://parking.duke.edu/buses/downtown-shuttle}.}
\end{itemize}

\subsection*{Course Objectives}
This course focuses on developing medical device prototyping skills that will be
used in future project and design courses, with a focus on preparing our
students for positions in the medtech industry.  Students will work individually
to complete design tasks that will be tested to quantitative specifications.
Students will gain hands-on experience with device fabrication, debugging,
testing and failure analysis.

Upon completion of this course, students should be able to:

\begin{itemize}
    \item Perform functional decomposition and develop user action flow chart
    \item Utilize ECAD (KiCad) for:
    \begin{itemize}
        \item Electronic schematic capture
        \item Printed circuit board (PCB) layout
        \item SPICE circuit simulation
    \end{itemize}
    \item Modular breadboarding of circuits to translation to testable PCBs
    \item Design a battery-powered device that:
    \begin{itemize}
        \item Accepts analog and digital user input
        \item Outputs analog and digital outputs
        \item Performs logic using discrete electronics (e.g., logic gates, 555 timers)
        \item Passes testable specifications to 95% CI
    \end{itemize}
    \item Implement electronic logic on microcontroller (ESP32 or equivalent platform)
    \begin{itemize}
        \item Modular / testable code development in C
        \item Software version control (\texttt{git})
        \item Interupt Service Routines
        \item Pulse Width Modulation
    \end{itemize}
    \item Utilize CAD (Onshape) for:
    \begin{itemize}
        \item Device enclosure design with UI/UX considerations
        \item Input / Output considerations
        \item Size / weight constraints
        \item Preparation of mechanical drawings
    \end{itemize}
    \item 3D printing of designs
    \begin{itemize}
        \item Assembly of discrete enclosure pieces
        \item Use of mechanical fasteners
    \end{itemize}
    \item Write technical analysis documents / reports
\end{itemize}


\subsection*{Prerequisites}
\begin{itemize}
    \item Introduction to Design (EGR101)
    \item Introductory Circuit Analysis (ECE110 or equivalent, corequisite)
\end{itemize}

\subsection*{Resources}
\subsection*{Hardware}
The following development kit is \textbf{required} hardware for this semester:
\textbf{TBD}

\subsubsection*{Textbooks}
All of these books contain valuable content and will be referenced throughout
the semester.  If you are looking to pursue a career in medical device design,
then it may be worth having one of these as a reference.

\begin{itemize}
    \item \href{https://find.library.duke.edu/catalog/DUKE009381868}{Practical Electronics for Inventors (Scherz \& Monk) [Fourth Edition]}\footnote{This is available online through the Duke Library.}
    \item The Art of Electronics (Horowitz \& Hill) [Third Edition]
    \item \href{http://www.ulrich-eppinger.net/}{Product Design and Development (Ulrich, Eppinger)}
    \item \href{https://www.crcpress.com/Design-of-Biomedical-Devices-and-Systems/King-Fries-Johnson/p/book/9781466569133}{Design of Biomedical Devices and Systems, King, Fries \& Johnson}
\end{itemize}

\subsubsection*{Online Resources}
We will be using the following [E]CAD packages:
\begin{itemize}
    \item \href{https://duke.onshape.com}{Onshape}\footnote{Note, this is a
    specific Duke instance of Onshape, not the public one.  You will be invited
    to join by Dr. Palmeri.}
    \footnote{If you know Fusion 360 or SolidWorks and would rather use those CAD
    packages, you are welcome to, but the class will only support Onshape.}
    \item \href{https://kicad.org/}{KiCad}
    \footnote{Please download and install the stable version (currently \texttt{v6.0.10}). Do not install the
    \texttt{v7.x} release candidate!}
    \footnote{If you know Altium Designer, you are welcome to use it instead, but
    the class will only support KiCad.}
\end{itemize}

These are cloud tools or cross-platform for installation on your personal computers.

\begin{itemize}
    \item \href{https://teams.microsoft.com}{Teams} - ``Slack-like'' tool for live
    chat, team discussion, asking questions and resource sharing
    \item \href{https://diagrams.net}{diagrams.net} - software / hardware flowcharts, functional decomposition
    \item \href{https://git-scm.com/}{Git} - version control software
    \item \href{https://gitlab.oit.duke.edu}{GitLab} - version control / project mangement
\end{itemize}

\subsection*{Attendance \& Participation}
Class participation, including lecture attendance, team meetings, and scheduled
lab time, contributes to your class grade. Not being able to participate in
class activities due to
 illness should be reported using the
\href{http://www.pratt.duke.edu/undergrad/policies/3531}{Short Term Illness
Form (STIF)} \textbf{before} the missed class activity.

\textbf{If you have a non-illness-related reason for not being able participate
in a class activity or meet an assignment submission deadline, please reach out
to Dr.\ Palmeri as soon as possible to discuss the situation.}

Students are responsible for obtaining missed lecture content from other
students in the class.  Lectures will be recorded via Panopto and available via
Sakai.

\subsection*{Class Schedule} 
The following table is an overview of activities this semester (always use
Gradescope for the latest information on assignment due dates).

\begin{calendar}{1/10/2023}{16} 
    \setlength{\calboxdepth}{.5in}
    \setlength{\calwidth}{1.0\textwidth} 
    \skipday % Monday (no class)
    \calday[Tuesday]{\classday}
    \skipday % Wednesday (no class)
    \calday[Thursday]{\classday}
    \skipday % Friday (no class)
    \skipday\skipday % weekend (no class)

    \Holiday{1/16/2023}{MLK Holiday}
    \Holiday{3/13/2023}{Spring Break}
    \Holiday{3/15/2023}{Spring Break}
    \Holiday{3/17/2023}{Spring Break}

    \caltexton{2}{Syllabus Review, Functional Decomposition \& User Action Flowcharts}
    \caltexton{2}{\textbf{Lab:} Software Installations}
	\caltexton{3}{CAD: Intro to Workflows}
	\caltexton{4}{CAD: Assemblies \& Mechanical Drawings}
	\caltexton{4}{\textbf{Lab:} CAD (Flange \& Bushing)}
	\caltexton{5}{CAD: Tips on more complicated parts}
	\caltexton{6}{ECAD: Introduction to Schematic Capture}
	\caltexton{6}{\textbf{Lab:} CAD (Hybrid III 6 y/o)}
	\caltexton{7}{ECAD: Schematic Capture Best Practices}
	\caltexton{8}{ECAD: Pin Designation, ERC, Spice Simulation}
	\caltexton{8}{\textbf{Lab:} KiCad Schematic Capture}
	\caltexton{9}{Electronics: Power, Voltage Regulation, Filters, [Light Emitting] Diodes}
	\caltexton{10}{Electronics: Transistors, 555 Timers, Buttons (Debounce)}
	\caltexton{10}{\textbf{Lab:} Circuit Design / Analysis}
	\caltexton{11}{Timer Project Overview, Functional Decomposition, User Action Flowchart}
	\caltexton{12}{Project Circuit Design \& Schematic Capture}
	\caltexton{12}{\textbf{Lab:} Project Circuit Design \& Schematic Capture}
	\caltexton{13}{ECAD: Introduction to PCB Layout}
	\caltexton{14}{ECAD: PCB Layout (Power Traces, Grounding, Test Points)}
	\caltexton{14}{\textbf{Lab:} Project PCB Layout}
	\caltexton{15}{3D Printing, Fasteners, Panel-Mounted Components}
	\caltexton{16}{Project: Physical Enclosure}
	\caltexton{16}{\textbf{Lab:} Project CAD}
	\caltexton{17}{Project Integration}
	\caltexton{18}{Project Integration}
	\caltexton{18}{\textbf{Lab:} Project Integration, 3D Printing, PCB Fabrication}
	\caltexton{19}{PCB Testing}
	\caltexton{20}{PCB Testing}
	\caltexton{20}{\textbf{Lab:} PCB Testing}
	\caltexton{21}{Introduction to Microcontrollers}
	\caltexton{22}{Microcontrollers: ISR, PWM, Timers}
	\caltexton{22}{\textbf{Lab:} Project Integration \& Testing}
	\caltexton{23}{Crashcourse: C}
	\caltexton{24}{Crashcourse: Software Version Control (\texttt{git})}
	\caltexton{24}{\textbf{Lab:} Using \texttt{git} and flashing MCU}
	\caltexton{25}{MCU Project Overview (Blinking and Brightness Modulation)} 
	\caltexton{26}{Timer \& MCU Projects}
	\caltexton{26}{\textbf{Lab:} Timer \& MCU Projects}
	\caltexton{27}{Timer \& MCU Projects}
	\caltexton{28}{Timer \& MCU Projects}
	\caltexton{28}{\textbf{Lab:} Timer \& MCU Projects}
	\caltexton{29}{LDOC: Summary of everything}

\end{calendar}

Final Project Due: May 01, 2023 at 09:00

\subsection*{Grading} 
The following grading scheme will be used for this course:

\begin{center}
\begin{tabular}{ll}
Participation                     & 15\% \\
Midterm Projects \& Deliverables  & 50\% \\
Final Project Report              & 35\% \\
\end{tabular}
\end{center}

\subsubsection*{Final Course Grades}
This course is not ``curved'' (i.e., a distribution of grades will not be
enforced), and a traditional grading scheme will be used (e.g., 90-93 = A-,
94-97 = A, 97-100 = A+).

Failing the course can happen with a cummulative score $<$ 65 or not completing
all of the assignments.

All grades will be posted to Gradescope throughout the semester to track your performance.

\subsubsection{Gradescope}
All graded assignments, associated due dates (usually Fridays at 17:00), and
grades/feedback will be posted on Gradescope.  Regrades must be submitted via
Gradescope.
\begin{itemize}
	\item \textbf{You must associate the pages of your submission with the
	grading rubric criteria during the submission process.}  Failure to do
	so will result in lost assignment credit.
	\item For team submissions, please make sure that assignments have
	\textbf{each team member associated with the submission}.
\end{itemize}

\subsubsection*{Late Policy}
Permission to submit an assignment late should be sought from Dr. Palmeri as
far in advance as reasonably possible, but no less than 48 hours in advance,
except in cases of illness.

Unexcused late assignments will lose 50\% of their potential point value for
each 24 hour period beyond the due date (i.e., 100\% \textrightarrow
 50\% → 25\%\ldots).

\textbf{All assignments must be satisfactorily completed by each student to pass
the class, even if no credit will be awarded based on the late policy.}

\subsubsection*{Regrades}
Any regrading requests need to be made via Gradescope within one week of grades for a given
assignment being returned using Gradescope. You
 must provide a description of
why you feel a regrade is appropriate. Requesting a
 regrade could lead to
additional loss of credit when an assignment is re-evaluated.

There will be a combination of individual and team-graded assignments. Some assignments will have an opportunity to be
resubmitted based on grading feedback at the discretion of Dr. Palmeri.

\subsection*{Duke Community Standard \& Academic Honor}
Engineering is inherently a collaborative field, and in this class, you are
encouraged to work collaboratively on your projects. The
 work that you submit
must be the product of your and your group’s effort and understanding. All
resources developed by another person or company, and used in your project,
must be properly acknowledged.

All students are expected to adhere to all principles of the
\href{https://trinity.duke.edu/undergraduate/academic-policies/community-standard-student-conduct}{Duke
Community Standard}. Violations of the Duke Community
 Standard will be referred
immediately to the Office of Student Conduct. Please do not hesitate to talk
with Dr. Palmeri about any
 situations involving academic honor, especially if
it is ambiguous what should be done.

\end{document}
